  \subsubsection{床関数の差で境界・条件・繰り上がりを検出する}\label{節：ガウス型の境界と条件}
    正整数$n$と整数$m\ge2$に対して,
    \[
      \left\lfloor\frac nm\right\rfloor-
      \left\lfloor\frac{n-1}{m}\right\rfloor
      =\begin{cases}
         1 & (n\text{が}m\text{の倍数}),\\
         0 & (\text{それ以外})
       \end{cases}
    \]
    となる.商の整数部分が増える境界を検出しているのである.
    同様に,\,$\lfloor\sqrt n\rfloor-\lfloor\sqrt{n-1}\rfloor$は$n$が平方数のときだけ1になる.
    \textbf{式の形で条件を表し,条件に応じて変化する\bookterm{recurrence}{漸化式}を作れる}という使い方である.

    二つの数の和でも,整数の境界を越えたかを調べられる.
    小数部分を$\{x\}=x-\lfloor x\rfloor$と書けば,
    \begin{align*}
      &\lfloor x+y\rfloor-\lfloor x\rfloor-\lfloor y\rfloor\\
      &\qquad=\lfloor\{x\}+\{y\}\rfloor\in\{0,1\}.
    \end{align*}
    これは小数部分の和で繰り上がりが起きたときだけ1になる.
    \bookterm{floor-function}{床関数}は一般には和に分配できないが,その差には意味がある.

    \noindent \bookblocklink{example-gauss-role-4}{problem}{answer}{【例題】}\par\nobreak
    \begin{enumerate}[label=(\arabic*),parsep=3pt,beginpenalty=10000,format=\bookitemlink{example-gauss-role-4}{problem}{answer}]
      \item \bookterm{sequence}{数列}を次のように定める.\bookterm{general}{一般項}を求め,\bookterm{difference}{階差}の意味を説明せよ.
            \[
              a_0=0,\qquad a_n=a_{n-1}+d_n\quad(n\ge1),
            \]
            \[
              \begin{aligned}
                d_n={}&\left\lfloor\frac n4\right\rfloor
                       -\left\lfloor\frac{n-1}{4}\right\rfloor\\
                      &-\left\lfloor\frac n6\right\rfloor
                       +\left\lfloor\frac{n-1}{6}\right\rfloor.
              \end{aligned}
            \]
    \end{enumerate}

    \noindent \bookblocklink{example-gauss-role-4}{hint}{problem}{【方針】}\par\nobreak
    \begin{enumerate}[label=(\arabic*),parsep=3pt,beginpenalty=10000,format=\bookitemlink{example-gauss-role-4}{hint}{problem}]
      \item 各差を倍数かどうかの判定として読みます.
            和を取ると,途中の\bookterm{floor-function}{床関数}が打ち消し合います.
    \end{enumerate}

    \noindent \bookblocklink{example-gauss-role-4}{answer}{problem}{【解答】}\par\nobreak
    \begin{enumerate}[label=(\arabic*),parsep=3pt,beginpenalty=10000,format=\bookitemlink{example-gauss-role-4}{answer}{problem}]
      \item \[
              a_n=\left\lfloor\frac n4\right\rfloor
                  -\left\lfloor\frac n6\right\rfloor.
            \]
            \,$d_n$は4の倍数だけなら1,6の倍数だけなら$-1$,両方またはいずれでもないなら0である.
    \end{enumerate}

    \noindent \bookblocklink{example-gauss-role-4}{solution}{problem}{【解法】}\par\nobreak
    \begin{enumerate}[label=(\arabic*),parsep=3pt,beginpenalty=10000,format=\bookitemlink{example-gauss-role-4}{solution}{problem}]
      \item \bookterm{floor-function}{床関数}の差の性質から,\,$d_n$は【解答】に述べた条件を表す.
            また,\,$a_n=\sum_{k=1}^{n}d_k$であり,
            \begin{align*}
              \sum_{k=1}^{n}\left(
                \left\lfloor\frac k4\right\rfloor-
                \left\lfloor\frac{k-1}{4}\right\rfloor\right)
                &=\left\lfloor\frac n4\right\rfloor,\\
              \sum_{k=1}^{n}\left(
                \left\lfloor\frac k6\right\rfloor-
                \left\lfloor\frac{k-1}{6}\right\rfloor\right)
                &=\left\lfloor\frac n6\right\rfloor.
            \end{align*}
            両式の差から\bookterm{general}{一般項}を得る.\,$n=0$でも両辺は0である.
            条件を一つずつ場合分けして値を追う方法と,差を足して一度に求める方法がある.
    \end{enumerate}
